% =====================================================================
%  Harald Høffding, Review of Claudius Wilkens, Erkendelsens Problem
%  Translated from the Danish: „Claudius Wilkens: Erkendelsens Problem (Anmeldelse)“, Det nittende Aarhundrede (October 1875--March 1876), October 1875, pp. 81--85.
%
%  TRANSLATION -- GENERATED; DO NOT EDIT BY HAND.  Built by ebook-method/hoeffding-1875-76/tr_parts.py build wilkens-anmeldelse-1875 --write
%  from the parts in ebook-method/hoeffding-1875-76/tr/ (GLOSSARY.md, TRANSLATOR-BRIEF.md).  Made from the collated
%  transcription texts/hoeffding/wilkens-anmeldelse-1875/transcription.tex (2026-10-15); [N] in the margin = the page of the original.
%  Terms: GLOSSARY.md (choices marked ◆ there for Hans to review).
% =====================================================================
\documentclass[12pt,a4paper,openany]{book}
\usepackage[utf8]{inputenc}
\usepackage[T1]{fontenc}
\usepackage{libertinus}
\usepackage{libertinust1math}
\usepackage[german,french,english]{babel}
\usepackage[protrusion=true,expansion=true]{microtype}
\usepackage{setspace}
\usepackage[margin=1.3in]{geometry}
\usepackage{fancyhdr}
\usepackage{hyperref}
\hypersetup{pdftitle={Review of Wilkens, Erkendelsens Problem}, pdfauthor={Harald Høffding}, hidelinks}
\pagestyle{fancy}
\fancyhf{}
\fancyhead[LE,RO]{\thepage}
\fancyhead[C]{\textit{H. Høffding: Review of Wilkens, Erkendelsens Problem (1875)}}
\fancypagestyle{plain}{\fancyhf{}\fancyhead[LE,RO]{\thepage}}
\setstretch{1.3}
\usepackage{marginnote}
\usepackage{graphicx}
\DeclareUnicodeCharacter{0254}{\reflectbox{c}}   % ɔ (id est)
\renewcommand{\thefootnote}{\ifcase\value{footnote}\or *\or **\or ***\else\arabic{footnote}\fi)}
\newcommand{\opage}[1]{\marginnote{\footnotesize\textit{[#1]}}}
\newcommand{\signature}[1]{\par\smallskip\hfill #1\hspace*{2em}\par}
\pdfstringdefDisableCommands{\renewcommand{\footnote}[2][]{}\renewcommand{\opage}[1]{}}

\begin{document}

\begin{titlepage}
\begin{center}
{\LARGE\bfseries The Nineteenth Century.}\\[0.6em]{\small A Monthly for Literature and Criticism}\\[0.4em]{\small edited by}\\[0.2em]{\large Georg Brandes and Edvard Brandes.}\\[0.8em]{\small 1875. October--December, and 1876. January--March.}\\\vfill{\small Copenhagen. Published by Gyldendal (F. Hegel). 1876.}\\[2em]
{\small Translated from the Danish}
\end{center}
\end{titlepage}
\clearpage

% --- p. 81 ---
\setcounter{footnote}{0}%
\opage{81}\begin{center}{\large \emph{Literature}.}\\[0.3em]\rule{3em}{0.4pt}\end{center}\medskip

{\small\noindent \emph{Claudius Wilkens}: \textit{Erkendelsens Problem. Et Forsøg paa positiv Kriticisme} [The Problem of Knowledge: An Attempt at a Positive Criticism]. Copenhagen 1875.}\medskip

The author of the work before us makes it very difficult, even for sympathetic readers, to grant him the approval his
work really deserves, because in the writing itself he has beforehand taken up and assigned to himself far more than even
those readers can concede him. It has done his work great harm. The reviewer knows several instances of readers who were
frightened off from making themselves acquainted with the book on account of the preface. Philosophical investigations
ought, especially in these skeptical times, and especially when one, like the author, expressly proclaims that one stands
on an exact standpoint, to appear with greater modesty; that is to the benefit both of science and of its devotees.

But this must not distract attention from the fact that the author is working in a direction and in a field that may be
said to be on the order of the day. Ever since the dissolution of the Hegelian school the conviction has been making its
way that idealist speculation actually regarded itself, wrongly, as the consummating consequence of Kant's great
philosophical reform. This speculation was a natural reaction against Kant's subjective and revolutionary standpoint, but
it has far from exhausted, or even merely had a full eye for, the significance of his doctrine. Recent natural science,
especially the physiology of the senses, and the realistic school in England have also awakened attention to the great
significance of Kant's fundamental thought in all its simplicity. That we do not apprehend things as they are in
themselves, but as they present themselves to us on account of our organization, on this there is now ever more
agreement, and no investigator has seen it as clearly as
% --- p. 82 ---
\setcounter{footnote}{0}%
\opage{82}Kant. It was the problem of the origin of knowledge that led Kant to this insight. Over against the realistic
doctrine that all knowledge, all content of consciousness, was impressed from without into consciousness, stood the
idealistic doctrine of innate ideas, which, if they did not lie ready in the soul from the first, yet gradually, as it
were, uncovered themselves from their wrappings. Kant settled the dispute by holding that we take nothing up into our
consciousness without its undergoing a transformation conditioned by our whole organization, but that on the other side
nothing is innate in us except this organization itself. Space, time and necessity and similar determinations therefore
belong to things only in and for our knowledge; but the intuition of space does not lie ready in us from the first; not
it, but the ground of it, is innate in us, just as the idea of the color red is not innate in us, whereas it has its
ground in our innate nature that we get the sensation of red when we see cinnabar.

The author calls his theory criticism, because like Kant he teaches that we do not know things in themselves. But he
calls it \emph{positive criticism}, because he believes he can fill up the gulf between knowledge and the thing in
itself. For though we do not know things in themselves, we do yet, in his opinion, know their relations, we have ``an
absolute knowledge of the relations of things''. The Kantian theory of knowledge is corrected by this, that a subjective
and an objective side is ascribed to the forms of knowledge as well as to its matter''. --- This last expression is
unfortunate. Kant ascribes ``an objective side'' to the forms of knowledge too; that we intuit things in space is,
according to him, just as much grounded in things as in our own nature; but we do not know what it is in things that
corresponds to our intuition of space. Since our forms of intuition and of thought spring from our nature, so that we can
never abstract from them but see everything under the presupposition of them, Kant had complete right to conclude that
we do not know what things in themselves are. Since space is a \emph{form of intuition}, it is self-contradictory to
assert that things \emph{in themselves, i.e. abstracted from every observer} are in space. Had the author known the
discussion carried on about this in Germany in recent years, he would hardly have spoken so confidently of Kant's
fallacy here. In particular Hermann Cohen, in a uniquely thorough and acute work (``Kants Theorie der Erfahrung'',
Berlin 1871), has shown Kant's consistent line of thought on this point.

% --- p. 83 ---
\setcounter{footnote}{0}%
\opage{83}The self-contradiction mentioned comes out clearly in one of the author's own examples. If we have two sticks, the one
twice as large as the other, then the relation between them --- says the author --- must surely be 1/2, whether we look
at them or not, thus in itself. But a relation exists only on the presupposition of a comparison; when there is no one
who compares, one can therefore not speak of the relation 1/2 either. That it has its ground in the essence of things
that we apprehend them in mathematical relations, Kant does not, as said, deny. But the author maintains an
\emph{absolute} knowledge of the relations of things; things themselves we apprehend, by his concession, only through
symbols, but the relation between the symbols and the relation between the things are supposed to be one and the same,
so that the author even wants to have it expressed by a mathematical equation. Kant could concede that our intuitions
correspond to the relations of things, but not that they resemble them, let alone that they are equal to them
(``gemäss'', not ``ähnlich'', let alone ``gleich'').

It is difficult to form a definite idea of how the author actually wants his doctrine of the absolute knowledge of the
relations of things to be understood. At the oral defense of the dissertation he declared it to be his opinion that the
relations between the ideas were the
same as (=) the relations between the things. But in the dissertation itself he expresses himself very vacillatingly.
Pag. 25 (cf. P. 118) it is said that we get indeed only a subjective picture of the world, but a correct picture of the
relations of existence as if knowledge could not be correct because it is subjective!
Is it not correct, e.g., that cinnabar is red, and yet this is a subjective sensation, for cinnabar is red only for our
eyes. --- Pag. 131, 177 and 182, on the other hand, call the apprehension of the relations of things symbolic, which is
the same as subjective; here the author has then forgotten his theory. --- Pag. 29 the author appeals to positivism, as
agreeing with his theory, and Pag. 117 it is said: ``a predictive knowledge is absolute''. Here the author confuses an
exact, i.e. perfectly distinct and quantitatively determinate knowledge with an absolute knowledge in the sense of one
that is not bound to the phenomenal but concerns that which is in itself. Exact science concerns only phenomena, but
gives us a completely distinct insight into their relations; positivism seeks to gather the results of all exact
sciences in one system; but neither does it
% --- p. 84 ---
\setcounter{footnote}{0}%
\opage{84}go beyond the phenomenal. --- This vacillation and obscurity is a natural consequence of the self-contradictory
character of the theory.

The author's ``correction of Kant'' is thus both superfluous and incorrect. But he believes he still has another door
open to the knowledge of the Absolute, of the thing in itself. ``Self-observation'', ``self-feeling'', ``the feeling of
life'' open for us a bridge to the essence of existence. Our own being we sense immediately and have thus here one point
of existence opened up for us; we know how we ourselves fare by being, and infer by the way of analogy from that how
other existing beings fare. This appeal to feeling in order to stop a hole in the theory does not augur well. The author
has himself sketched this point so indistinctly that it has surely only been his intention to feel his way forward. What
ambiguity lies not in a sentence like this: ``The deepest understanding is in reality that of feeling'' (Pag.
140)! The author here takes ``understanding'', as one of his opponents at the oral disputation pointed out, in a quite
different sense than elsewhere in the dissertation; to understand is now no longer to comprehend in a rational manner,
but to sympathize in a poetic way: in this sense one may well say with the author that we understand the animal better
than the stone; but of a concept there is then no longer any question. It is an illusion to think that we should better
comprehend the mental than the bodily, because the former is given to us immediately in our consciousness. To
comprehending belongs seeing the matter in its relation to everything else, but then precisely that immediacy is a
hindrance. The psychological phenomena are just as much phenomena as the bodily ones; and there exists no special logic
for them.

Pag. 123 the author defines ``feeling one's being'' by ``apprehending the worth of one's being'' and declares that a being
that neither possesses nor is itself the object of such a feeling is ``worthless, meaningless, i.e. senseless''. Here we
are referred to the practical field. Morally and historically we are indeed interested only in that which has ``value and
meaning''; but what vouches for us that this is something Absolute? The highest standard that we can lay at the
foundation of valuation is after all only relative, drawn from the demands and tasks of human nature, and it is a bold
dogmatism that the author here plunges into by identifying these with the core of existence, a dogmatism of a similar
kind to that into which Kant's ``moral faith'' led astray so many Kantians
% --- p. 85 ---
\setcounter{footnote}{0}%
\opage{85}and which the young Schelling combated in his interesting ``Briefe über Dogmatismus und Kriticismus''. The
sentence just cited is found almost in the same terms in Lotze (Mikrokosmus I, Pag. 407). But it is worth noting with
what caution this clear and acute thinker proceeds on this point. It is not in the theory of knowledge that he utters
this sentence; it is not even in metaphysics. When metaphysics is to have scientific character as a rational development
of the fundamental determination of our view of the world, it cannot take up practical, aesthetic-ethical postulates.
Therefore, according to Lotze (Metafysik Pag. 135), ``the expression the true being, which is only of aesthetic
significance [i.e. expresses a determination of value] is to be excluded from metaphysical usage''. Such circumspection
the author does not observe, although he fights (or wants to fight) under the banner of criticism.

There are certainly points enough where Kant's theory must be corrected. Recent English psychology has opened a series
of elementary investigations which partly rectify, partly carry through, what Kant wanted. Recent natural science and
inductive logic have brought forward entirely new problems. But on the whole and in the large the old Königsberger will
certainly for a long time yet stand with unassailed greatness as the central thinker of modern times. His ideas form a
firm core, around which everything that the later development in the various scientific fields has brought to light will
be able to gather.
\signature{\textbf{Harald Høffding.}}

\end{document}
